\section{Excluding bounded algebraic types}\label{sec:bounded}

Other than the high alternating and classical cases, simple socles of bounded order can be discarded immediately by the available interval length: \(M=HS\) at an almost simple current parent with socle \(S\), so \(|S|\ge|M:H|\) tends uniformly above any prescribed constant. By CFSG the remaining cases use a finite list of root systems, with parameter as large as needed, \(S=(G^{F_A})'\) for a simple adjoint algebraic group over \(k=\overline{\mathbf F}_p\).

Throughout this section we work at a current almost simple parent in the prepared reverse-flat interval of \Cref{sec:patterns}. Thus $H\le M\le\operatorname{Aut}(S)$, $S$ is the simple socle, and $M=HS$. The root systems under consideration belong to a fixed finite list. Every complexity bound below depends only on this list and on the stated degree bound; it does not depend on the characteristic or the finite-field parameter.

\subsection{Uniform algebraic input}

\begin{definition}[Bounded equation complexity]\label{w:complexity}
Fix closed affine split models for the root systems in the finite list, with a bounded number of coordinates and bounded-degree group operations. A closed subset has equation complexity at most $d$ if it is the common zero set, inside the chosen group, of polynomials of degree at most $d$. There need not be a previously specified bound on the number of equations: their linear span lies in the finite-dimensional space of polynomials of degree at most $d$.
\end{definition}

\begin{theorem}[Sufficiently general finite subgroups]\label{w:lp}
For the fixed finite list of simple adjoint types, there is a uniform degree bound $d_0$ with the following property. If a finite subgroup $\Gamma\le G(k)$, where $k=\overline{\mathbf F}_p$, is contained in no proper closed subset of equation complexity at most $d_0$, then there is a Steinberg map $\Psi$ such that
\[
 (G^\Psi)'\le\Gamma\le G^\Psi.
\]
Here the degree bound follows from the uniform constructible-family
formulation of the sufficiently-general subgroup theorem of
Larsen--Pink \citep[Theorem~0.5 and Metadefinition~2.2]{LarsenPink2011}:
the exceptional family has uniformly bounded defining equations in
the fixed affine models.
\end{theorem}

This formulation applies to the fixed split models in all characteristics. Its hypothesis is sufficient generality, expressed by avoidance of bounded-complexity proper closed subsets. The alternative argument in \Cref{sec:quantitative} will also establish the implication needed over the indicated algebraic closures.

\begin{convention}[Point maps and height]\label{w:point-maps}
We also use the standard classification of the Steinberg maps and of automorphisms and orders of the finite Lie type groups. On points take the group of maps generated by inner and diagram automorphisms, split Frobenius powers and inverses, and special isogeny powers where applicable (square split \(p\)-Frobenius in special small characteristics). It has a height homomorphism to \(\mathbf Z\) with unit of parameter growth \(p\), or \(\sqrt p\) in the special case. Height zero maps are ordinary algebraic automorphisms, and positive height maps are Steinberg maps with fixed groups, of parameter \(P=(\text{unit})^{\rm height}\). The normalizer of \(S\) here induces the automorphisms required for our large-parameter \(S\). Modulo inners one uses the standard pinned maps here: split Frobenius commutes with diagram maps, and in the cases with the special self-isogeny one uses its powers (square as above, no other nontrivial ordinary diagram needed in those types). Thus arbitrary words simplify to an inner twist of these combinations; the positive-height ones have standard fixed-group types (inner twists removed by Lang--Steinberg). The exceptional generators only concern fixed small characteristics. All coordinates, variety complexities etc. may use fixed closed affine split models of bounded size/degree, with the group operations and pinned ordinary diagram maps of bounded complexity uniformly in the characteristic. Images of bounded-degree-equation varieties by the whole point-map group have bounded equation degree (field powers just transform coefficients, inner and diagram have bounded complexity, special map of fixed degree has inverse using itself and inverse field power). Positive maps are given algebraically of degree \(O(P)\). The classification and fixed-group facts used here are those of \citet{Steinberg1968,GLS1994}. We write $\mathcal A(G)$ for this group of maps on $G(k)$, $\operatorname{ht}$ for its height, and $\delta$ for the parameter unit, so $\delta=p$ or $\sqrt p$.
\end{convention}

\begin{lemma}[Order, degree, and pointwise rigidity]\label{w:height-rigidity}
For the finite type list, let $b=\dim G$ and $B(\Psi)=(G^\Psi)'$. For sufficiently large parameter $P=\delta^{\operatorname{ht}(\Psi)}$, the following hold.
\begin{enumerate}
\item $B(\Psi)$ is simple, and its order exceeds $|G^\Theta|$ for every $\Theta\in\mathcal A(G)$ with $0<\operatorname{ht}(\Theta)<\operatorname{ht}(\Psi)$.
\item For each fixed equation complexity $d$, $B(\Psi)$ lies in no proper closed subset of complexity at most $d$, once $P$ is sufficiently large for $d$.
\item If a positive-height point map fixes $B(\Psi)$ pointwise, it is a positive integral power of $\Psi$.
\end{enumerate}
\end{lemma}
\begin{proof}
Two useful uniform comparisons follow. Write \(b=\dim G\). For \(P\) large, \(B(\Psi)=(G^\Psi)'\) is simple of order at least \(P^b(1-o(1))/z\) inside \(G^\Psi\) of order \(P^b(1+o(1))\), with index at most a center bound \(z\) from the root system. In particular \(|B(\Psi)|\) exceeds the entire fixed-group order at any smaller positive height, since \(z<(\text{unit})^b\) by the usual center sizes. This applies uniformly to comparison with smaller heights: below parameter \(\sqrt P\) even the rough order bounds suffice, and otherwise the two product formulas both have small relative error. Also \(B(\Psi)\) lies in no proper closed subvariety of bounded degree, for \(P\) sufficiently large for that bound. On any irreducible component of dimension \(e<b\) and bounded degree, the fixed-point equations of degree \(O(P)\) cut a finite set of size \(O(P^e)\), by affine B\'ezout (combine the equations successively generically to cut to dimension zero). Hence a point map of positive height fixing \(B(\Psi)\) pointwise must be a power of \(\Psi\): divide its height by the given height and compose off that power on points. Positive remainder is impossible by orders, zero-height nonidentity by the bounded equations of its fixed locus.

For the last assertion, if $\alpha$ has height $q\operatorname{ht}(\Psi)+r$ with $0\le r<\operatorname{ht}(\Psi)$, the point map $\Psi^{-q}\alpha$ still fixes $B(\Psi)$ pointwise. If $r>0$, its fixed group is too small. If $r=0$, a nonidentity algebraic automorphism has a proper fixed locus with uniformly bounded equations, contradicting the second assertion. Hence $r=0$ and $\alpha=\Psi^q$.
\end{proof}

\subsection{Witness stabilizers and the divisor contradiction}

\begin{lemma}[A bounded-complexity failure witness]\label{w:witness}
Let $H\le K<M$ be a tested overgroup with $\Gamma=K\cap S\ne1$, and suppose $\Gamma$ fails the sufficient-generality hypothesis of \Cref{w:lp}. For sufficiently large parameter of $S$, there is a proper full stabilizer label $K'$ with $K\le K'<M$ such that
\[
 K'\cap S=N_G(Z)\cap S
\]
for a proper closed subset $Z\subset G$ whose equation complexity, and that of its conjugation normalizer $N_G(Z)$, are uniformly bounded.
\end{lemma}
\begin{proof}
Consider a proper \(K\) in the interval, above a tested large-index label, and \(\Gamma=K\cap S\ne1\). If insufficiently general for the sandwich, it lies in a proper bounded-degree-equation closed \(X\subset G\) by Larsen--Pink as used above. Take the subgroup of our point-map group preserving \(S\) and inducing \(M\)-elements on \(S\); it maps onto \(M\) by the automorphism description. Let \(\tilde K\) be the full preimage of \(K\) there (so each such map normalizes \(\Gamma\)). Then
\[
 Z=\bigcap_{\alpha\in\tilde K,\ \gamma\in\Gamma} \alpha(\gamma X)
\]
contains \(\Gamma\), is proper closed with equations of uniformly bounded degree, left-\(\Gamma\)-invariant and \(\tilde K\)-invariant. Even for the whole intersection one can just span all the equations occurring in the fixed bounded-degree space. Its point-map stabilizer (within that preimage of \(M\)) thus defines a label \(K'\ge K\), since the stabilizer before projection contains the kernel there (already in \(\tilde K\)). Identifying \(S\) with its inner actions on itself, \(K'\cap S=N_G(Z)\cap S\) by using the actual inner lifts to \(G\). The conjugation normalizer here has bounded-degree equations (test the equations of \(Z\) on all conjugates of its points by \(g,g^{-1}\)). It is proper since otherwise all conjugates of \(\Gamma\) left-stabilize \(Z\), while their normal closure is Zariski-dense by simplicity and trivial center. The degree comparison for \(S\) thus makes \(K'\) proper.

Here $N_G(Z)$ denotes the conjugation stabilizer, not the left stabilizer. To justify its properness explicitly, if every $g\in G(k)$ normalizes $Z$, left invariance under $\Gamma$ implies left invariance under every $g\Gamma g^{-1}$. The Zariski closure of the group generated by these conjugates is a nontrivial closed normal subgroup of the simple adjoint group $G$, hence is $G$. Since $Z$ is closed and nonempty, left invariance by this dense subgroup forces $Z=G$, a contradiction. The nonidentity of $\Gamma$ suffices because the adjoint group has trivial center.
\end{proof}

\begin{proposition}[No bounded-type parent with a stage remaining]\label{w:bounded-exclusion}
Fix the finite list of bounded root systems. There is a fixed trial increment $t$ such that a prepared current parent with a remaining private-chart stage cannot have socle of one of these types, provided the residual interval length is sufficiently large. All choices are uniform in the characteristic and parameter.
\end{proposition}
\begin{proof}
At the current parent flat $W$, test a sufficiently large fixed
increment $t$ in a private chart. We first find a flat $U$ of that
increment for which every $K\in[K_U,M]$ admits the sandwich for
$K\cap S$. If every trial failed, \Cref{w:witness} would supply full
stabilizer witnesses above the trial labels. The affine-pattern
construction gives an independent parameter list whose successive
label intersections decrease strictly for more steps than the dimension
of the relevant polynomial space. Descent remains strict after
intersection with $S$, since $M=HS$ and every overgroup of $H$ is
determined by its intersection with $S$. Each strict descent must add
a defining polynomial outside the span of the preceding equations.
But all conjugation-normalizer equations lie in one bounded-dimensional
space, a contradiction. This step needs only a bounded independent
list, no prior anchor, and tests relative to the current parent.

Put $E=K_U$ and $\Gamma_K=K\cap S$ for $E\le K\le M$. We now have
\[
 B(\Psi_K)\le\Gamma_K\le G^{\Psi_K}.
\]
The reserved interval length below $E$ makes every parameter here
large enough for \Cref{w:height-rigidity}. Indeed, if that residual
height is $h$, then
$|\Gamma_K|\ge|\Gamma_E|\ge[E:H]\ge2^h$, whereas all fixed groups
below a prescribed parameter have bounded order for the fixed type
list. Write $B_E=B(\Psi_E)$. Perfection and the abelian quotient of a
fixed group by its derived subgroup give $\Gamma_E'=B_E$.

Since $B_E\le\Gamma_K$, the map $\Psi_K$ fixes $B_E$ pointwise.
Rigidity therefore gives $\Psi_K=\Psi_E^{a_K}$ for a positive integer
$a_K$. The ambient map $F_A$ likewise fixes $B_E$, so $F_A=\Psi_E^d$.
It also fixes $B(\Psi_K)$ pointwise, and rigidity applied to that
simple group makes $F_A$ a positive power of $\Psi_K$. Consequently
\[
 F_A=\Psi_E^d,\qquad \Psi_K=\Psi_E^{a_K},\qquad a_K\mid d.
\]
These are equalities of point maps, so each exponent determines an
actual embedded fixed group.

Every lift of $E$ normalizes $\Gamma_E$ and hence $B_E=\Gamma_E'$.
Conjugating $\Psi_E$ by such a lift gives a map of the same height
fixing $B_E$ pointwise; rigidity makes it equal to $\Psi_E$. Thus the
lifts commute with $\Psi_E$. For each divisor $a$ of $d$, put
$B_a=B(\Psi_E^a)$. This group lies in $S$: its elements are fixed by
$F_A$, and perfection puts it in $(G^{F_A})'$. It is normalized by
$E$, so $E B_a$ is an overgroup label. Moreover
\[
 (E B_a)\cap S=\Gamma_E B_a\subseteq G^{\Psi_E^a},
 \qquad E B_1=E,\qquad E B_d=M.
\]
If $a\mid b$ with $a<b$, then $B_a\le B_b$, and
\Cref{w:height-rigidity} gives $|B_b|>|G^{\Psi_E^a}|$.
Hence $E B_a<E B_b$. A saturated divisor chain therefore gives
$\Omega(d)$ strict steps in the exact height-$t$ interval $[E,M]$,
where $\Omega(d)$ counts prime factors with multiplicity. Thus
\[
 \tau(d)=\prod_{\ell^a\parallel d}(a+1)
 \le2^{\Omega(d)}\le2^t.
\]

For a fixed exponent $a_K$, the sandwich leaves only boundedly many
choices for $\Gamma_K$: its quotient by $B_{a_K}$ is a subgroup of
the bounded-order group $G^{\Psi_E^{a_K}}/B_{a_K}$. If $z$ bounds
that order, $C=2^z$ bounds the number of choices. Since
$K=H\Gamma_K$, no additional choices of outer automorphisms occur.
We conclude that $|[E,M]|\le C2^t$.

The reverse-flat interval, however, contains at least
$q^{\lfloor(t-1)^2/4\rfloor}$ intermediate subflats, by the ordinary
linear-subspace count; this bound also covers an empty starting flat.
For sufficiently large fixed $t$ it exceeds $C2^t$, already using
$q\ge5$. This is the required contradiction, uniform in characteristic
and parameter, whenever a bounded-type parent has a stage remaining.
\end{proof}
